% p.37 — الباب الموفي بالعشرين في الافعال المعتلة وهي اثنتان ومائة
\mspage{37}
\bab{الباب الموفي بالعشرين في الافعال المعتلة\\{\normalsize وهي اثنتان ومائة}}
% Layout: six numbered columns (circled numbers above the heading row).
% Each column: category heading, باب of the verb, then ten rows (ماضي معروف،
% مضارع معروف، مصدر، اسم فاعل، ماضي مجهول، مضارع مجهول، مصدر، اسم مفعول، أمر،
% نهي) giving the uncontracted/original forms; below a rule the same ten rows
% give the forms after تعليل. × marks a non-existent form. In col. 1 the last
% two cells each hold two forms side by side (اُوْخُذْ خُذْ / لَا تَاخُذْ لَا تَخُذْ).
% The lithograph writes initial hamza as bare alif + vowel (اَخَذَ); the dagger
% alif it uses for the madda (اٰخِذٌ) is rendered آ.
% TODO: col. 6 lower block, rows 8–10 are smudged in the scan.
\begin{sarf}{|c|c|c|c|c|c|}
\thead{(1) مهموز الفاء} & \thead{(2) مهموز الفاء} & \thead{(3) مهموز الفاء} & \thead{(4) مهموز الفاء} & \thead{(5) مهموز الفاء} & \thead{(6) مهموز العين}\\
\hline
نَصَرَ & فَتَحَ & ضَرَبَ & كَرُمَ & سَمِعَ & فَتَحَ\\
\hline
اَخَذَ & اَهَبَ & اَدَبَ & اَدُبَ & اَرِجَ & سَأَلَ\\
يَأْخُذُ & يَأْهَبُ & يَأْدِبُ & يَأْدُبُ & يَأْرَجُ & يَسْأَلُ\\
اَخْذًا & اَهْبًا & اَدْبًا & اَدْبًا & اَرَجًا & سُؤَالًا\\
آخِذٌ & آهِبٌ & آدِبٌ & اَدِيْبٌ & اَرِجٌ & سَائِلٌ\\
اُخِذَ & اُهِبَ & اُدِبَ & \x & \x & سُئِلَ\\
يُؤْخَذُ & يُؤْهَبُ & يُؤْدَبُ & \x & \x & يُسْأَلُ\\
اَخْذًا & اَهْبًا & اَدْبًا & \x & \x & سُؤَالًا\\
مَأْخُوْذٌ & مَأْهُوْبٌ & مَأْدُوْبٌ & \x & \x & مَسْئُوْلٌ\\
اُؤْخُذْ & اِئْهَبْ & اِئْدِبْ & اُؤْدُبْ & اِئْرَجْ & اِسْأَلْ\\
لَا تَأْخُذْ & لَا تَأْهَبْ & لَا تَأْدِبْ & لَا تَأْدُبْ & لَا تَأْرَجْ & لَا تَسْأَلْ\\
\hline
اَخَذَ & اَهَبَ & اَدَبَ & اَدُبَ & اَرِجَ & سَأَلَ\\
يَاخُذُ & يَاهَبُ & يَادِبُ & يَادُبُ & يَارَجُ & يَسْأَلُ\\
اَخْذًا & اَهْبًا & اَدْبًا & اَدْبًا & اَرَجًا & سُوَالًا\\
آخِذٌ & آهِبٌ & آدِبٌ & اَدِيْبٌ & اَرِجٌ & سَائِلٌ\\
اُخِذَ & اُهِبَ & اُدِبَ & \x & \x & سُئِلَ\\
يُوْخَذُ & يُوْهَبُ & يُوْدَبُ & \x & \x & يُسْأَلُ\\
اَخْذًا & اَهْبًا & اَدْبًا & \x & \x & سُوَالًا\\
مَاخُوْذٌ & مَاهُوْبٌ & مَادُوْبٌ & \x & \x & \unclear{مَسُوْلٌ}\\
اُوْخُذْ خُذْ & اِيْهَبْ & اِيْدِبْ & اُوْدُبْ & اِيْرَجْ & \unclear{اِسْأَلْ سَلْ}\\
لَا تَاخُذْ لَا تَخُذْ & لَا تَاهَبْ & لَا تَادِبْ & لَا تَادُبْ & لَا تَارَجْ & \unclear{لَا تَسْأَلْ}\\
\end{sarf}
